\documentclass[a4paper, 12pt]{article}
\usepackage[english, russian]{babel}
\usepackage[T2A]{fontenc}
\usepackage[utf8]{inputenc}
\usepackage{textcomp}
\usepackage{amsmath}
\usepackage{amssymb}
\usepackage{amsfonts}

\usepackage{graphicx}
\usepackage{import}
\usepackage{pdfpages}
\usepackage{transparent}
\usepackage{xcolor}

% programming language support
\usepackage{listings}
\usepackage{color}

\definecolor{dkgreen}{rgb}{0,0.6,0}
\definecolor{gray}{rgb}{0.5,0.5,0.5}
\definecolor{mauve}{rgb}{0.58,0,0.82}

\lstset{frame=tb,
  language=C++,
  aboveskip=3mm,
  belowskip=3mm,
  showstringspaces=false,
  columns=flexible,
  basicstyle={\small\ttfamily},
  numbers=none,
  numberstyle=\tiny\color{gray},
  keywordstyle=\color{blue},
  commentstyle=\color{dkgreen},
  stringstyle=\color{mauve},
  breaklines=true,
  breakatwhitespace=true,
  tabsize=4
}

\graphicspath{ {figures/} }

\newcommand{\incfig}[2][1]{%
    \def\svgwidth{#1\columnwidth}
    \import{./figures/}{#2.pdf_tex}
}

\pdfsuppresswarningpagegroup=1

\usepackage{mdframed}
\usepackage{lipsum}

\usepackage{fancyhdr}
\pagestyle{fancy}
\fancyhf{} % Очистить текущие настройки колонтитулов

% Верхний колонтитул
\fancyhead[L]{\today}
\fancyhead[C]{\thepage}
\fancyhead[R]{Artem Svirin}
\renewcommand{\headrulewidth}{1pt} % Толщина линии

% Нижний колонтитул
\fancyfoot[L]{\today}
\fancyfoot[C]{\thepage}
\fancyfoot[R]{Artem Svirin}
\renewcommand{\footrulewidth}{1pt} % Толщина линии

\newmdtheoremenv{theo}{Теорема}
\newmdtheoremenv{proof_theo}{Доказательство теоремы} % Доказательство теоремы
\newmdtheoremenv{cor}{Следствие}[theo] % Следствие
\newmdtheoremenv{proof_cor}{Доказательство следствия}[theo] % Доказательство следствия
\newmdtheoremenv{zm}{Замечание}
\newmdtheoremenv{mdef}{Определение}
\newmdtheoremenv{stm}{Утверждение}

\title{Colloquium}
\author{Artem Svirin}
\date{\today}

\begin{document}
\begin{titlepage}
	\maketitle
	\thispagestyle{empty}
\end{titlepage}
\clearpage
\begin{enumerate}
	\item Определение числового ряда и его сходимости. Необходимое условие сходимости числового ряда. \\
	      \begin{mdframed}
		      \textbf{Определение числового ряда:} \\
		      Пусть дана последовательность $a_k$. Символ \\
		      \[
			      \sum_{k=1}^{\infty} a_k = a_1 + a_2 + \ldots + a_k
		      \]
		      Называется числовым рядом с общим членом $a_k$. \\
		      \textbf{Сходимость числового ряда:} \\
		      Ряд с общим членом $a_k$ называется сходящимся, если $\exists \lim_{n \to \infty}S_n$ в $\mathbb{R}$ \\
		      \textbf{Необходимое условие сходимости числового ряда:} \\
		      Если ряд $\sum_{k=1}^{\infty} a_k$ сходится, то $\lim_{k \to \infty} a_k = 0$
	      \end{mdframed}
	\item Критерий Коши сходимости ряда. Определение абсолютной и условной сходимости. \\
	      \begin{mdframed}
		      \textbf{Критерий Коши сходимости ряда:} \\
		      Ряд $\sum_{k=1}^{\infty} a_k$ сходится, тогда и только тогда, когда: \\
		      \[
			      \forall \varepsilon > 0 \ \exists n_0 = n_0(\varepsilon) \in \mathbb N: \ \forall n > n_0, \ \forall p \in \mathbb{N} \ \Rightarrow \ \left| \sum\limits_{k = n+1}^{n+p} a_k \right| < \varepsilon.
		      \]
		      \textbf{Определение абсолютной сходимости} \\
		      Если ряд $\sum |a_k|$ - сходится, то говорят, что ряд $\sum a_k$ сходится абсолютно. \\
		      \textbf{Определение условной сходимости} \\
		      Если ряд $\sum a_k$ сходится, а ряд $\sum |a_k|$ расходится, то говорят, что ряд $\sum a_k$ сходится условно.
	      \end{mdframed}
	\item Интегральный признак Коши. \\
	      \begin{mdframed}
		      \textbf{Интегральный признак Коши:} \\
		      Если $f(x) \ge 0$ невозрастающая функция на $[k; + \infty]$,  $k$ - натуральное, то  $\sum_{n=k}^{\infty} f(n)$ и $\int_{k}^{\infty}f(x)dx$ сходится или расходится одновременно. \\
		      \textbf{Более простая формулировка:} \\
		      Рассмотрим положительный числовой ряд $\sum_{n=1}^{\infty} a_n$. Если существует несобственный интеграл $\int_{1}^{\infty} a_n dx$, то ряд сходится или расходится вместе с этим интегралом.
	      \end{mdframed}
	\item Признаки сравнения в обычной и предельной форме. \\
	      \begin{mdframed}
		      \textbf{Признак сравнения в обычной форме:} \\
		      Если $0 \le a_n \le b_n$ с некоторого номера $n > N$, то из сходимости $\sum b_n \implies$ сходимость $\sum a_n$. \\
		      из расходимости $\sum a_n \implies$ расходимость $\sum b_n$. \\
		      \textbf{Признак сравнения в предельной форме:} \\
		      Пусть $a_n \ge 0, b_n \ge  0$, $\lim_{n \to \infty} \frac{a_n}{b_n} = k, k \neq 0, \infty$, тогда $\sum_{n=1}^{\infty} a_n$ и $\sum_{n=1}^{\infty} b_n$ сходятся или расходятся одновременно.
	      \end{mdframed}
	\item Признак Даламбера в обычной и предельной форме. \\
	      \begin{mdframed}
		      \textbf{Признак Даламбера в обычной форме:} \\
		      Пусть с некоторого номера $n, a_n > 0$ и  $\frac{a_{n+1}}{a_n} < q, q < 1$, тогда ряд $\sum_{n=1}^{\infty} a_n$ сходится.
		      Если $\frac{a_{n+1}}{a_n} \ge 1$, то $\sum_{n=1}^{\infty} a_n$ расходится. \\
		      \textbf{Признак Даламбера в предельной форме:} \\
		      Пусть с некоторого номера $n$,  $a_n > 0$ и существует \\  $\lim_{n \to \infty} \frac{a_{n+1}}{a_n} = q$, тогда:
		      \begin{itemize}
			      \item Если $q < 1$, то ряд $\sum_{n=1}^{\infty} a_n$ сходится.
			      \item Если $q > 1$, то ряд $\sum_{n=1}^{\infty} a_n$ расходится.
			      \item Если $q = 1$, то признак не работает. (требуется другой признак)
		      \end{itemize}
	      \end{mdframed}
	\item Радикальный признак Коши в обычной и предельной форме. \\
	      \begin{mdframed}
		      \textbf{Радикальный признак Коши в обычной форме:} \\
		      Пусть с некоторого номера $n, a_n >  0$ и  $\sqrt[n]{a_n} < q, q < 1$, тогда ряд $\sum_{n=1}^{\infty} a_n$ сходится.
		      Если $\sqrt[n]{a_n} \ge 1$, то $\sum_{n=1}^{\infty} a_n$ расходится. \\
		      \textbf{Радикальный признак Коши в предельной форме:} \\
		      Пусть с некоторого номера $n$,  $a_n > 0$ и существует \\  $\lim_{n \to \infty} \sqrt[n]{a_n} = q$, тогда:
		      \begin{itemize}
			      \item Если $q < 1$, то ряд $\sum_{n=1}^{\infty} a_n$ сходится.
			      \item Если $q > 1$, то ряд $\sum_{n=1}^{\infty} a_n$ расходится.
			      \item Если $q = 1$, то признак не работает. (требуется другой признак)
		      \end{itemize}
	      \end{mdframed}
	\item Определение знакочередующегося ряда. Ряд Лейбница. Теорема Лейбница. \\
	      \begin{mdframed}
		      \textbf{Определение знакочередующегося ряда:} \\
		      ряд $\sum_{n=1}^{\infty} (-1)^{n}b_n, b_n > 0$ или $\sum_{n=1}^{\infty} (-1)^{n+1}b_n, b_n > 0$ называется знакочередующимся. \\
		      \textbf{Ряд Лейбница:} \\
		      Ряд Лейбница - это знакочередующийся ряд, у которого\\ $\lim_{n \to \infty} b_n = 0$ и $b_{n+1} \le b_n$ \\
		      \textbf{Теорема Лейбница:} \\
		      Всякий ряд Лейбница сходится (вообще говоря, условно) и его остаток по модулю не превосходит модуля первого отброшенного члена.
	      \end{mdframed}
	\item Перестановка ряда. Теорема Римана. \\
	      \begin{mdframed}
		      \textbf{Перестановка ряда:} \\
		      Перестановкой $n$ называется биекция $\sigma: \mathbb{N} \to \mathbb{N}$ \\
		      Перестановкой ряда $\sum_{n=1}^{\infty} a_n$ называют ряд $\sum_{k=1}^{\infty} a_{\sigma(k)}$ \\
		      \textbf{Теорема Римана:} \\
		      Если ряд $\sum_{n=1}^{\infty} a_n$ сходится условно, то существует перестановка ряда, сход. к любому выбранному числу,  к $\pm \infty$.
	      \end{mdframed}
	\item Перестановка абсолютно сходящегося ряда. \\
	      \begin{mdframed}
		      \textbf{Перестановка абсолютно сходящегося ряда:} \\
		      Если ряд $\sum_{n=1}^{\infty} a_n$ сходится абсолютно, то любая его перестановка сходится к тому же числу.
	      \end{mdframed}
	\item Произведение рядов. \\
	      \begin{mdframed}
		      \textbf{Произведение рядов:} \\
		      Пусть даны два ряда $\sum_{n=1}^{\infty} a_n$ и $\sum_{n=1}^{\infty} b_n$.  $\sum_{k=1}^{\infty} a_{m(k)}b_{n(k)}$ \\
		      пара $(m(k), n(k))$ - счетное множество, существует \\ биекция $k \in \mathbb{N} \to (m(k), n(k))$  \\

	      \end{mdframed}
	\item Признак Дирихле и Абеля: \\
	      \begin{mdframed}
		      \textbf{Признак Дирихле:} \\
		      ряд $\sum_{k=1}^{\infty} a_k b_k$ сходится, если:
		      \begin{itemize}
			      \item Последовательность частичных сумм $A_k = \sum_{l=1}^{k} a_l$ ограничена. (т.е $\exists M: \forall k \ \Rightarrow \ |A_k| \le M$)
			      \item Последовательность $\{b_k\}^{\infty}_{k=1}$ монотонна и стремится к 0.
		      \end{itemize}
		      \textbf{Признак Абеля:} \\
		      ряд $\sum_{k=1}^{\infty} a_k b_k$ сходится, если:
		      \begin{itemize}
			      \item последовательность $\{b_k\}^{\infty}_{k=1}$ монотонна и ограничена.
			      \item $\sum_{k=1}^{\infty} a_k$ - сходится.
		      \end{itemize}
	      \end{mdframed}
	\item Признаки Раабе и Гаусса. \\
	      \begin{mdframed}
		      \textbf{Признак Раабе:} \\
		      $\sum p_n$ сходится, если  $\forall n>n_0$ \\
		      $\frac{p_n}{p_{n+1}} \le  1- \frac{\alpha}{n}$, где $\alpha > 1$ и \\
		      $\sum p_n$ расходится, если с некоторого номера $\frac{p_{n+1}}{p_n} \ge 1-\frac{1}{n}$ \\
		      \textbf{Признак Раабе в предельной форме:} \\
		      Если $p_n > 0$, и существует $\lim_{n \to \infty} n \left( \frac{p_n}{p_{n+1}} - 1 \right) = l$, тогда:
		      \begin{itemize}
			      \item Если $l > 1$, то ряд $\sum_{n=1}^{\infty} p_n$ сходится.
			      \item Если $l < 1$, то ряд $\sum_{n=1}^{\infty} p_n$ расходится.
			      \item Если $l = 1$, то признак не работает. (требуется другой признак)
		      \end{itemize}
		      \textbf{Признак Гаусса:} \\
		      $a_n > 0, \varepsilon > 0$ и $\frac{a_n}{a_{n+1}} = \lambda + \frac{M}{n}+O \left( \frac{1}{n} + \varepsilon \right)$
		      \begin{itemize}
			      \item Если $\lambda > 1$, то ряд $\sum_{n=1}^{\infty} a_n$ сходится.
			      \item Если $\lambda < 1$, то ряд $\sum_{n=1}^{\infty} a_n$ расходится.
			      \item Если $\lambda = 1$, то:
			            \begin{itemize}
				            \item Если $M > 1$, то ряд $\sum_{n=1}^{\infty} a_n$ сходится.
				            \item Если $M \le  1$, то ряд $\sum_{n=1}^{\infty} a_n$ расходится.
			            \end{itemize}
		      \end{itemize}
	      \end{mdframed}
	\item Методы суммирования рядов. Формула суммирования Эйлера и Абеля. \\
	\item Гипергеометрическое дифференциальное уравнение. Гипергеометрическая функция Гаусса. Формула суммирования в единице. \\
	\item Определение функционального ряда и его сходимости, а также области сходимости. \\
	      \begin{mdframed}
		      \textbf{Определение функционального ряда:} \\
		      Пусть дана последовательность функций $f_k(x)$. \\
		      \[
			      \sum_{k=1}^{\infty} f_k(x) = f_1(x) + f_2(x) + \ldots + f_k(x)
		      \]
		      Называется функциональным рядом с общим членом $f_k(x)$. \\
		      \textbf{Сходимость функционального ряда в точке $x_0$:} \\
		      $\lim_{n \to \infty} S_n(x_0) = S(x_0)$, \\
		      $\forall \varepsilon > 0 \exists N = N(\varepsilon): \forall n > N(\varepsilon) \Rightarrow |S_n(x_0) - S(x_0)| < \varepsilon$ \\
		      \textbf{Область сходимости функционального ряда:} \\
		      Областью сходимости функционального ряда называется множество точек $x$, в которых ряд сходится.
	      \end{mdframed}
	\item Определение равномерной сходимости функционального ряда. Критерий Коши и теоерема Вейерштрасса. \\
	      \begin{mdframed}
		      \textbf{Определение равномерной сходимости функционального ряда:} \\
		      Функциональный ряд $\sum_{k=1}^{\infty} f_k(x)$ сходится равномерно на множестве $E$, если: \\
		      $\forall \varepsilon > 0 \ \exists N = N(\varepsilon): \forall n > N, \forall x \in E \Rightarrow |S_n(x) - S(x)| < \varepsilon$ \\
		      \textbf{Критерий Коши равномерной сходимости функционального ряда:} \\
		      Функциональный ряд $\sum_{k=1}^{\infty} f_k(x)$ сходится равномерно на множестве $E$, если: \\
		      $\forall \varepsilon > 0 \ \exists N = N(\varepsilon): \forall n > N, \forall p \in \mathbb{N}, \forall x \in E \Rightarrow |S_{n+p}(x) - S_n(x)| < \varepsilon$ \\
		      \textbf{Теорема Вейерштрасса:} \\
		      Если $|f_k(x)| \le a_k, \forall x \in E, \forall k \in \mathbb{N}$ и $\sum_{k=1}^{\infty} a_k$ сходится, то ряд $\sum_{k=1}^{\infty} f_k(x)$ сходится равномерно на $E$.
	      \end{mdframed}
	\item Свойство почленного интегрирования функционального ряда. \\
	      \begin{mdframed}
		      \textbf{Свойство почленного интегрирования функционального ряда:} \\
		      $\sum_{n=1}^{\infty} u_n(x) \underset{[a,b]}\rightrightarrows u_n(x), u_n(x) \in R[a,b] \implies$ \\
		      $\implies \int_{{a}}^{{b}} {\sum_{n=1}^{\infty} u_n(x)} \: d{x} {} = \sum_{n=1}^{\infty} \int_{{a}}^{{b}} {u_n(x)} \: d{x}$
	      \end{mdframed}
	\item Свойство почленного дифференцирования функционального ряда. \\
	      \begin{mdframed}
		      \textbf{Свойство почленного дифференцирования функционального ряда:} \\
		      $\sum_{n=1}^{\infty} u_n(x)$ можно почленно дифференцировать, если $u'_n(x) \in C^{1}([a;b])$, \\
		      $\sum_{n=1}^{\infty} u_n(x_0)$ существует ряд сходящийся в нашей точке $x_0$ и $\sum_{n=1}^{\infty} u'_n(x) \underset{[a;b]}\rightrightarrows$. \\
		      Тогда $\left( \sum_{n=1}^{\infty} u_n(x) \right)' = \sum_{n=1}^{\infty} u'_n(x)$
	      \end{mdframed}
	\item Определение степенного ряда. Теорема Абеля.
	      \begin{mdframed}
		      \textbf{Определение степенного ряда:} \\
		      Степенным рядом называется ряд вида: \\
		      \[
			      \sum_{n=0}^{\infty} c_nx^n = c_0 + c_1x + c_2x^2 + \ldots + c_nx^n + \ldots
		      \]
		      \[
			      \sum_{n=0}^{\infty} c_n(x-a)^n = c_0 + c_1(x-a) + c_2(x-a)^2 + \ldots + c_n(x-a)^n + \ldots
		      \]
		      $a = const$ \\
		      \textbf{Теорема Абеля:} \\
		      Если степенной ряд $\sum_{n=0}^{\infty} a_nx^n$ сходится в некоторой точке $c>0$, то он сходится на $[0;c]$ \\
		      Если есть сходимость в точке  $c, (c>0)$, то ряд сходится на  $[-c;0]$ и  $S(x)$ непрерывно на  $[-c;0]$ \\
	      \end{mdframed}
	\item Теорема о существовании радиуса сходимости. Определение радиуса сходимости и формулы его нахождения. \\
	      \begin{mdframed}
		      \textbf{Теорема о существовании радиуса сходимости:} \\
		      $ \sum_{n=0}^{\infty} c_n z^n (2)$ \\
		      Для всякого степенного ряда $(2)$ существует  $R (R\ge 0$ - число или $+\infty)$, такое что: \\
		      \begin{itemize}
			      \item Если $R \neq 0$ и $R \neq +\infty$, то ряд абсолютно сходится в круге $K = \{z: |z|< R\}$ и расходится вне круга $K$. Этот круг называют кругом сходимости, а $R$ - радиусом сходимости. \\
			      \item Если $R = 0$, то ряд сходится только в точке $z = 0$. \\
			      \item Если $R = +\infty$, то ряд сходится во всей комплексной плоскости. \\
		      \end{itemize}
		      \textbf{Формулы нахождения радиуса сходимости:}
		      \begin{itemize}
			      \item $R = \lim_{n \to \infty} \left| \frac{a_n}{a_{n+1}} \right|$
			      \item $R = \lim_{n \to \infty} \frac{1}{\sqrt[n]{|a_n|}}$
		      \end{itemize}
	      \end{mdframed}
\end{enumerate}


\end{document}
